\section{Correctness and Cost}\label{sec:correct}

The engine has two obligations that do not follow from the standard argument for
a subsumption-based search: it must terminate although convex merging destroys
the antichain property, and its verdict must be the correct one although its
stored set depends on the schedule.  We discharge both from a single property of
the update, which also justifies the two-phase protocol, and then bound the cost
of each mechanism of \cref{sec:approach}.

\subsection{The contract and the probe}

\begin{proposition}\label{prop:contract}
\textsc{Update} of \cref{alg:update} satisfies \emph{(U1)--(U3)}, for every
merge window $m \ge 0$ and whether or not the stored set has a mirror.
\end{proposition}
\begin{proof}
\emph{(U1).}  The procedure returns \textsc{dropped} without mutating anything
exactly when line~2 or fold~(a) fires.  Fold~(a) tests $\Zone \subq_q Y$ over
every live stored zone, so it fires whenever $\mathrm{cov}_\Passed(q,\Zone)$
holds; line~2 fires only when $\Zone \subseteq Y$ for some live $Y$, which
implies $\Zone \subq_q Y$ by \cref{lem:preorder}, so adding it neither creates
nor removes drops.  Conversely a drop exhibits a stored $Y$ with $\Zone \subq_q
Y$, so coverage held.

\emph{(U2).}  Let $R$ be the set of zones retired by folds~(b) and~(c) and
$\Zone^{*}$ the zone finally interned.  Fold~(c) replaces the current candidate
$Y_t$ by $\hull(Y_t,Y'_t)$, which equals $Y_t \cup Y'_t$ because the merge
condition requires the union to be exact, and retires $Y'_t$; after $r$ merges
$\Zone^{*} = \Zone \cup Y'_1 \cup \dots \cup Y'_r$.  Hence $\Zone \subseteq
\Zone^{*} \subseteq \Zone \cup \bigcup R$, and each $Y'_t \subseteq \Zone^{*}$
gives $Y'_t \subq_q \Zone^{*}$ by \cref{lem:preorder}.  A zone $Y$ retired by
fold~(b) satisfies $Y \subq_q \Zone \subseteq \Zone^{*}$, so $Y \subq_q
\Zone^{*}$ by \cref{lem:preorder} and transitivity.  Both arguments are
independent of how many merges were attempted, so they hold for every $m$; and
they read stored zones only, which is what a mirror changes.

\emph{(U3).}  Every call at $q$ holds the exclusive lock of $\stripe(q)$, and
$\stripe$ is a function of $q$.\qed
\end{proof}

\begin{lemma}[Monotonicity of coverage]\label{lem:monotone}
Let $\Passed$ and $\Passed'$ be the stored sets at $q$ before and after one
update.  Then $\mathrm{cov}_\Passed(q,X)$ implies $\mathrm{cov}_{\Passed'}(q,X)$
for every zone $X$.  Consequently, along every interleaving, once a symbolic
state is covered it stays covered.
\end{lemma}
\begin{proof}
Let $Y \in \Passed(q)$ with $X \subq_q Y$.  Under (U1) the stored set is
unchanged.  Under (U2), either $Y \notin R$, so $Y \in \Passed'(q)$; or $Y \in
R$, so $Y \subq_q \Zone^{*} \in \Passed'(q)$ and transitivity
(\cref{lem:preorder}) gives $X \subq_q \Zone^{*}$.  The second statement follows
by induction over the sequence of updates at $q$, which is well defined by
(U3).\qed
\end{proof}

\Cref{lem:monotone} does not say that the stored set grows: zones are retired and
its size can fall.  It says that the set of covered symbolic states grows, which
is the only thing a coverage probe reports on.

\begin{theorem}[The probe is exact for the case it answers]\label{thm:probe}
If a thread's probe at $(q,\Zone)$ returns \emph{true} and the thread skips the
update, the resulting passed set and the set of states recorded as inserted are
identical to those obtained by running the update at any later point of the same
interleaving.
\end{theorem}
\begin{proof}
The probe holds the shared lock of $\stripe(q)$, which excludes every writer of
that stripe, so it reads one well-defined stored set $\Passed_t(q)$, and
returning \emph{true} means $\mathrm{cov}_{\Passed_t}(q,\Zone)$.  For any later
$t'$, \cref{lem:monotone} gives $\mathrm{cov}_{\Passed_{t'}}(q,\Zone)$, so an
update executed at $t'$ takes branch (U1): it leaves the stored set unchanged and
records nothing, which is what skipping does.\qed
\end{proof}

\Cref{thm:probe} separates this protocol from optimistic concurrency, which reads
without exclusion, acts, and re-validates because its observation may have been
invalidated.  Here a positive observation cannot be invalidated, so there is
nothing to re-validate and the mutating phase is not entered at all.

\subsection{Termination and verdict}

\begin{theorem}\label{thm:term}
Assume the candidates offered at each discrete state $q$ fall into finitely many
classes of $\Zone \equiv_q \Zone' :\iff \abs_q(\Zone) = \abs_q(\Zone')$.  Then
the engine performs finitely many insertions and terminates, for every
interleaving.
\end{theorem}
\begin{proof}
Fix $q$ and let $\Zone_1,\Zone_2,\dots$ be the candidates inserted at $q$ in the
total order given by (U3).  Suppose $\Zone_i \equiv_q \Zone_j$ for some $i < j$.
Then $\Zone_j \subseteq \abs_q(\Zone_j) = \abs_q(\Zone_i)$, i.e.\ $\Zone_j
\subq_q \Zone_i$.  Immediately after the insertion of $\Zone_i$ the stored set
contains $\Zone_i^{*} \supseteq \Zone_i$, so $\mathrm{cov}(q,\Zone_i)$ held then,
and by \cref{lem:monotone} it still holds when $\Zone_j$ is offered: there is a
stored $Y$ with $\Zone_i \subq_q Y$, hence $\Zone_j \subq_q Y$ by transitivity,
so the update of $\Zone_j$ takes branch (U1) and does not insert --- contradicting
the choice of $\Zone_j$.  The inserted candidates therefore lie in pairwise
distinct classes, and their number is bounded by the number of classes at $q$.
Summing over the finite set $\dstates$ bounds the total number of insertions.
The frontier of level $k+1$ consists of zones inserted during level $k$, so a
level that inserts nothing ends the search; the number of levels is therefore
finite and each expands a finite frontier.\qed
\end{proof}

The hypothesis holds in both subsumption modes: under \cref{eq:alu} by
\cref{lem:finite}, and under inclusion because every candidate is then
$\extraLU$-extrapolated and those form a finite set~\cite{Daws98,Behrmann06}.
The proof uses neither an antichain invariant nor any order on updates beyond
(U3).  This is what makes it survive fold~(c), which can leave a stored $Y$ with
$Y \subq_q \Zone^{*}$ un-retired, and parallel execution, which fixes no order
across discrete states.

\begin{lemma}[Exact unions are transparent to $\succfn$]\label{lem:union}
If $Y = \Zone_1 \cup \Zone_2$ as sets of valuations, then for every edge, or
synchronised pair of edges, of $\net$ the zone produced from $Y$ is the union of
the zones produced from $\Zone_1$ and from $\Zone_2$.
\end{lemma}
\begin{proof}
The successor computation composes conjunction with a guard or invariant, clock
reset, and time elapse.  Each is defined pointwise on valuations --- $\{v : v
\models g\} \cap S$, $\{v[x{:=}0] : v \in S\}$ and $\{v + d : v \in S, d \ge 0\}$
--- and every operation of that shape commutes with union.\qed
\end{proof}

\begin{proposition}\label{prop:sound}
At every point of every interleaving, every $Y \in \Passed(q)$ satisfies $Y
\subseteq \bigcup\{\Zone : (q,\Zone) \text{ reachable}\}$.
\end{proposition}
\begin{proof}
Induction over the updates in the order given by (U3).  Initially $\Passed(q_0) =
\{\Zone_0\}$.  Branch (U1) changes nothing.  Branch (U2) adds $\Zone^{*}
\subseteq \Zone \cup \bigcup R$, where $\Zone$ is the zone of a generated
symbolic state and each $Y \in R$ satisfies the claim by hypothesis; by
\cref{lem:union} and the hypothesis, every generated zone at $q$ is contained in
the union of the reachable zones at $q$.\qed
\end{proof}

\begin{theorem}\label{thm:verdict}
For every interleaving, the engine returns \emph{true} if and only if some
reachable symbolic state satisfies $\varphi$.
\end{theorem}
\begin{proof}
\emph{Soundness.}  The engine returns \emph{true} only after evaluating $\varphi$
on a generated symbolic state $(q,\Zone_g)$; stored zones are never offered to
$\varphi$.  By \cref{lem:union,prop:sound}, $\Zone_g$ lies in the union of the
zones of reachable symbolic states at $q$.  As $\varphi$ is a condition on $q$
conjoined with a clock constraint tested for non-empty intersection, a witness
inside that union is a witness inside one of its operands.

\emph{Completeness.}  A zone inserted during level $j$ is stamped $j$ and, if
still live at the barrier of level $j$, is collected into $\Frontier_{j+1}$ and
expanded during level $j+1$; so at termination every live stored zone has been
expanded, the final frontier being empty.  Let $(q,\Zone)$ be reachable with
$\varphi(q,\Zone)$; we induct on the length $d$ of a path to it.  For $d = 0$ the
initial state is tested explicitly.  For $d > 0$ write $(q,\Zone) \in
\succfn(q_p,\Zone_p)$.  A second induction on $d$ shows
$\mathrm{cov}_\Passed(q_p,\Zone_p)$ at termination: the candidate covering
$(q_p,\Zone_p)$ was either dropped, in which case coverage held then and still
holds by \cref{lem:monotone}, or inserted, in which case coverage held
immediately after.  Let $Y_p$ be a witness; it was expanded, and \cref{lem:sim}
gives $(q,\Zone_2) \in \succfn(q_p,Y_p)$ with $\varphi(q,\Zone_2)$.  That
successor was generated and tested.\qed
\end{proof}

\Cref{thm:verdict} is the precise sense in which the engine is deterministic.
The stored set is not: it depends on the order in which candidates reach a
discrete state.  The verdict, and the set of covered symbolic states at
termination, do not.

\begin{proposition}[The abstraction switch is sound mid-run]\label{prop:adaptive}
Enabling $\extraLU$ at a barrier preserves \emph{(U1)--(U3)}, hence
\cref{thm:term,thm:verdict}.
\end{proposition}
\begin{proof}
The switch changes which zone is offered as a candidate and nothing about the
passed-set discipline: the same preorder $\subq_q$ is used before and after, and
the contract refers only to it, so a stored set mixing exact and widened zones is
covered by the same argument as a homogeneous one.  In \cref{thm:verdict},
completeness uses $\Zone \subq_q \Zone_2 \subseteq \extraLU(\Zone_2)$ and
transitivity; soundness needs $\extraLU$ to preserve $\varphi$, which holds
because $L_q$ and $U_q$ are floored by the clock constants of $\varphi$, so
$\varphi$ is preserved for the same reason a guard
is~\cite{Behrmann06,Daws98}.  Termination is unaffected, the switch being enabled
only when $\subq_q$ is \cref{eq:alu}, whose finite index already discharges the
hypothesis of \cref{thm:term}.\qed
\end{proof}

\subsection{What the update's cost mechanisms preserve}\label{sec:cost}

\begin{lemma}[Inclusion against a minimal constraint system]\label{lem:minclosure}
Let $\mg$ be a set of difference constraints with $\Zone' = \mathrm{close}(\mg)$
non-empty and let $\Zone$ be canonical.  Then $\Zone \subseteq \Zone'$ if and
only if $\Zone[i][j] \le b$ for every $(i,j,b) \in \mg$.
\end{lemma}
\begin{proof}
($\Rightarrow$)  Closure only tightens, so $\Zone'[i][j] \le b$ for every
$(i,j,b) \in \mg$; with \cref{eq:denseincl}, $\Zone[i][j] \le \Zone'[i][j] \le
b$.

($\Leftarrow$)  Fix $i,j$.  If no path from $i$ to $j$ exists in the constraint
graph of $\mg$ then $\Zone'[i][j] = \infty$.  Otherwise $\Zone'[i][j]$ is the
weight of a shortest such path $i = v_0 \to \dots \to v_r = j$ with edge bounds
$b_1,\dots,b_r$, so $\Zone'[i][j] = \sum_t b_t$.  As $\Zone$ is canonical,
repeated use of the triangle inequality gives $\Zone[i][j] \le \sum_t
\Zone[v_{t-1}][v_t]$; by hypothesis $\Zone[v_{t-1}][v_t] \le b_t$, and addition
of bounds is monotone, so $\Zone[i][j] \le \sum_t b_t = \Zone'[i][j]$.  With
$i,j$ arbitrary, \cref{eq:denseincl} yields the claim.\qed
\end{proof}

\Cref{lem:minclosure} replaces an $O(n^3)$ closure followed by an $O(n^2)$
comparison with an $O(|\mg|)$ scan of the packed record, needs no scratch matrix,
and is exact rather than an approximation.  The engine prunes with $\subq_q$,
which under \cref{eq:alu} reads arbitrary entries of the stored zone and is
therefore not decidable from $\mg$; the two are composed instead, the exact test
first as a filter.

\begin{theorem}[The filter preserves the explored state set]\label{thm:samestates}
The update with the packed-record filter drops exactly the candidates the
update without it drops, and stores exactly the zones it stores.
\end{theorem}
\begin{proof}
If line~2 fires on a live $Y$ then $\Zone \subseteq Y$ by
\cref{lem:minclosure}, hence $\Zone \subq_q Y$ by \cref{lem:preorder}, so
fold~(a) would also have dropped $\Zone$; both outcomes leave the stored set
unchanged, so which zone witnessed the drop is immaterial.  If line~2 fires on
no zone, folds (a)--(c) run unchanged.\qed
\end{proof}

\Cref{thm:samestates} is what makes the filter safe to enable unconditionally: it
is a cheaper decision procedure for a test the engine already performs, not a new
abstraction with its own soundness burden.  It is also the reason the coverage
probe is incomplete on a stored set without a mirror --- inclusion is strictly
finer than \cref{eq:alu}, so a candidate that is subsumed but not included falls
through to the exclusive phase, at the cost of one lock.

\begin{proposition}[The mirror is result-invariant]\label{prop:mirror}
The update reads the same zones in the same order with and without a mirror.
Consequently a mirror changes neither the stored set nor the verdict, and
refusing one for want of budget is likewise result-invariant.
\end{proposition}
\begin{proof}
It suffices to establish the invariant that, whenever the stripe lock is not
held, reading the live slots of a mirror in reverse order yields exactly the
sequence of live zones obtained by walking the chain of $q$ from its head.  It
holds at creation, which copies the live zones in reverse chain order, the chain
being prepended to.  It is preserved by the four operations that touch a mirror,
all under the exclusive lock: insertion appends the zone that is simultaneously
prepended to the chain; retirement clears the mirror slot in the same critical
section that marks the chain node; rebuilding removes cleared slots and preserves
the relative order of the rest; and the barrier compaction, which renumbers chain
nodes, rewrites the mirror's references to them.  Given the invariant, folds
(a)--(c) enumerate the same operands in the same order, and they read stored
zones while mutating only the candidate.  A refused mirror leaves the chain walk
in place, which by the same invariant enumerates the same zones.\qed
\end{proof}

\begin{proposition}[Work-based promotion bounds the wasted unpacks]
\label{prop:mirrorwork}
Let a stored set be promoted at its $v^{*}$-th visit under the rule $v \cdot L
\ge W$.  The unpacks performed on it before promotion number at most $W
\mathcal{H}_{v^{*}-1} = O(W \log v^{*})$, with $\mathcal{H}$ the harmonic number.
Under a rule $L \ge H$ on size alone, a stored set with $L < H$ is never promoted
and performs $\Theta(v \cdot L)$ unpacks, unbounded in $v$.
\end{proposition}
\begin{proof}
Write $L_v$ for the size at the $v$-th visit, which unpacks $L_v$ zones.  For
$v < v^{*}$ the rule has not fired, so $L_v < W/v$, and $\sum_{v<v^{*}} L_v < W
\sum_{v<v^{*}} 1/v$.  The second claim is immediate, the rule inspecting only
$L$.\qed
\end{proof}

The rule generalises the size rule rather than replacing it: a stored set that
already holds $W$ live zones is promoted on its first visit.  Rebuilding a mirror
is likewise amortised: it runs only once half the slots are cleared and costs
$O(L)$, which charged to the $\Omega(L)$ retirements that triggered it is $O(1)$
per retirement.

\begin{proposition}[The merge window is unconstrained]\label{prop:window}
For every $m \ge 0$ the engine terminates and returns the same verdict.
\end{proposition}
\begin{proof}
Immediate from \cref{prop:contract}, which holds for every $m$, together with
\cref{thm:term,thm:verdict}, which use the update only through its contract.\qed
\end{proof}

\Cref{prop:window} says $m$ may be chosen on cost grounds alone, not that the
choice is without consequence: merging changes what is stored, so $m$ trades
precision against time in both directions.  A narrow window misses the merges
that collapse a stored set and stores more zones; an unbounded window pays a
quadratic number of merge attempts and, merging being non-confluent, need not
store fewer.

\subsection{The cost of a barrier}\label{sec:barriercost}

\begin{proposition}[Bucketed deduplication is exact]\label{prop:bucket}
Let $b$ be any function of the discrete state.  Deduplicating $\{q \in R_t : b(q)
= j\}$ over all threads $t$ separately, for each bucket $j$, yields exactly
$\bigcup_t R_t$ and scans each discrete state once.
\end{proposition}
\begin{proof}
Every occurrence of $q$, in whichever thread's record, is placed in bucket
$b(q)$, since $b$ depends on $q$ alone.  That bucket observes all occurrences of
$q$ and no other bucket observes any.\qed
\end{proof}

\begin{proposition}[Compaction is amortised constant]\label{prop:compact}
Rebuilding a stripe once its retired nodes reach a third of its nodes costs
$O(1)$ amortised per retirement.
\end{proposition}
\begin{proof}
Let $N$ be the stripe's nodes and $D$ its retired nodes when the rebuild fires,
so $3D \ge N$.  The rebuild walks every node once, at cost $O(N) = O(D)$, and
leaves no retired node, so its cost is charged to the $D$ retirements that
enabled it, each charged once.\qed
\end{proof}

\begin{proposition}[Amortised reclamation]\label{prop:gcamort}
Reclamation performs $O(1)$ work per interned record, amortised, and a run
performs $O(\log)$ reclamations in the size of the live set.
\end{proposition}
\begin{proof}
A reclamation triggered at allocation level $A$ costs $O(A)$; afterwards
$\mathit{next} = 2A_{\mathrm{live}}$, so the next cannot occur before a further
$A_{\mathrm{live}}$ records have been interned.  Charging the $O(A)$ cost to
those $\Omega(A_{\mathrm{live}})$ interns, and using $A \le 2A_{\mathrm{live}}$
at the trigger, gives $O(1)$ amortised per intern; the trigger doubles, so the
number of reclamations is logarithmic in the final live set.\qed
\end{proof}

\begin{proposition}[Bounded overhead]\label{prop:gcmem}
Between reclamations, unreclaimed garbage never exceeds the live set, so the zone
table stays within twice its pinned-to-live size.
\end{proposition}
\begin{proof}
A reclamation fires as soon as $A$ reaches $2A_{\mathrm{live}}$, and $A -
A_{\mathrm{live}}$ bounds the garbage.\qed
\end{proof}

\Cref{prop:gcmem} is the bound one wants before making reclamation lazier.  An
absolute threshold on the high-water index count differs in exactly the way that
matters: that counter never falls, so it latches at the first large level and
forces a mark and sweep at every barrier thereafter, a serial cost growing with
the state space.

\begin{proposition}[Index recycling forces memo invalidation]\label{prop:gcsound}
If the zone table recycles freed indices in place and does not erase freed
records, then every memo mapping record content to an index, or an index to a
matrix, must be invalidated at each reclamation.
\end{proposition}
\begin{proof}
Reclamation frees only indices unreferenced by a live chain node or a frontier
entry, so live references stay valid.  A content-keyed memo, however, validates a
hit by comparing record bytes, and a freed record still compares equal to the
zone it held.  A thread interning that zone after the reclamation is handed a
freed index, stores it in a new chain node, and the next intern popping that
index from the free list overwrites the record; the node then keeps its discrete
state and denotes a different zone, so the search can reach symbolic states the
network cannot.  An index-keyed memo fails for the same reason in the other
direction.\qed
\end{proof}
